\subsection{معرفی و انگیزه}
این مکانیزم بر مبنای مبانی نظری پیشرفته در ادبیات توجه و فشرده‌سازی \cite{vaswani2017attention} توسعه یافته است.

مکانیزم توجه با سرهای پرس‌وجوی گروه‌بندی‌شده (\lr{Grouped-Query Attention - GQA}) به عنوان استانداردی موثر برای کاهش بار حافظه بافر کلید/مقدار (\lr{KV-Cache}) در مدل‌های ترنسفورمر به کار می‌رود. با این وجود، در روش‌های مرسوم \lr{GQA}، نسبت انتساب سرهای پرس‌وجو به سرهای کلید/مقدار به صورت یکنواخت و صلب در تمامی لایه‌ها و موقعیت‌های متنی یکسان در نظر گرفته می‌شود. این رویکرد ایستا بدون توجه به اهمیت محتوایی و انتروپی اطلاعاتی توکن‌ها عمل کرده و ممکن است به افت کیفیت بازنمایی در سرهای بحرانی منجر گردد.

مکانیزم توجه آگاه از کیفیت و ظرفیت (\lr{Quality and Capacity-Aware GQA - QCQA}) با هدف پویاسازی و یادگیری توزیع تخصیص سرها معرفی شده است. ایده محوری ماژول \lr{QCQA}، ترکیب پیوسته و هوشمند سرهای کلید/مقدار فشرده‌شده از طریق یک شبکه مسیریاب هموار (\lr{Continuous Learned Routing Distribution}) با دمای کنترل‌شده است. بدین ترتیب، برای هر سر پرس‌وجو، یک بازنمایی مجازی سنتز‌شده از کلید و مقدار بر مبنای ترکیبی خطی و وزن‌دار از سرهای فشرده موجود ساخته می‌شود که بدون نیاز به تکثیر صلب داده‌ها، انعطاف بیان‌پذیری مدل را به حداکثر می‌رساند.
